\section{Introduction}\label{sec:introduction}

Evaluating an option strategy requires a model of how contract values change as market conditions evolve. A position that appears profitable at expiration may incur substantial losses before that date, when a trader may need to close or hedge it. Stress tests and trading-system studies therefore require paths of option values and risk sensitivities as well as paths of the underlying stock price. This requirement is particularly important for American options because their values also reflect the opportunity to exercise before expiration. Historical option quotes provide observations of these quantities over realized market conditions, but a simulation must supply them for stock paths that have not occurred. Implied volatility (IV) provides a useful connection between market quotes and an option-pricing model because it is the volatility input recovered by inverting an observed price~\cite{black1973}. To use this connection in a synthetic market, the simulator needs a rule for assigning IV across strikes and maturities and updating those values as the stock path unfolds. The stock model must also represent the return behavior relevant to the intended experiment, including features such as heavy tails and volatility clustering~\cite{cont2001}. The resulting modeling problem is to produce complete option-contract paths whose response to these stock scenarios can be examined and interpreted.

Existing approaches address different parts of this problem through their choices of return dynamics, volatility representation, and pricing model. Parametric models such as stochastic alpha beta rho (SABR) and stochastic volatility inspired (SVI) provide compact descriptions of volatility smiles~\cite{hagan2002,gatheral2004}. Neural methods allow more flexible representations and can accelerate pricing and calibration within specified volatility models~\cite{horvath2021,bayer2019}. Models of IV-surface dynamics describe changes in level, skew, and curvature through common factors~\cite{cont2002}, with more recent methods incorporating static arbitrage constraints into surface simulation~\cite{cont2023}. Local- and stochastic-volatility models instead place stock and volatility dynamics within a joint pricing model~\cite{dupire1994,derman1994,heston1993}. These approaches already support generation beyond the observed quotes. A different design question arises when the researcher wants to retain a physical stock-return model chosen for its empirical behavior and vary the volatility assumptions separately. Historical and generative return models can reproduce important features of equity data~\cite{cont2001,wiese2020}, but their stock paths alone do not specify an option-implied volatility surface. For example, the Jump Hidden Markov Model (JumpHMM) generates returns through Markov transitions among return states and explicit visits to tail states~\cite{jumphmm2025}. Those returns determine terminal option payoffs while leaving the IV input for interim valuation unspecified. An evolving contract-level IV model can supply that input to a recombining American lattice~\cite{cox1979,leisen1996}. This separation allows researchers to compare volatility assumptions on the same stock paths, provided that the accuracy and consistency of the resulting option prices are evaluated explicitly.

In this study, we developed a computational framework that connected JumpHMM stock paths to fitted IV surfaces, contract-specific square-root factors, and American-option valuation. We fitted the surfaces to observed option data before simulation and used them to set variance targets as moneyness and remaining time to expiration changed. Each contract's variance factor evolved around its target and supplied the volatility input for a lattice calculation of option value and price sensitivities. This construction allowed the IV dynamics to change without replacing the physical return generator or using future option quotes during simulation. We evaluated the surface representation on 234{,}549 observations from 31 tickers and fifteen snapshots, comparing global, sector-specific, and per-ticker parameter sharing. Temporal holdouts examined transfer to withheld dates and the contribution of earnings-calendar inputs. We then simulated Goldman Sachs (GS) and Eli Lilly (LLY) contracts and conducted a paired ablation that held stock paths fixed across five IV variants. This experiment measured how the IV assumptions affected interim marks and early-liquidation profit and loss (P\&L), with identical terminal payoffs providing a control for interpretation. Lattice refinement and strike checks assessed whether the measured differences reflected numerical error or inconsistencies across contracts. We extended the surface evaluation across 53 chronological folds and fitted a separate July-cutoff model for comparison with later GS and LLY option observations. These tests distinguished conditional surface transfer, joint stock--option forecasting, and a diagnostic supplied with the subsequently observed stock path. To examine the stock contribution separately, we compared the retained JumpHMM with unchanged-price, adaptive-volatility, and simple directional benchmarks in 2025 and August--September 2026, selecting the new models' settings using earlier data. This evaluation tested the accuracy of the stock inputs without changing the original IV ablation. The resulting evidence addressed the framework's use for controlled comparisons and the limits of its forecasting performance, while joint no-arbitrage consistency remained a separate requirement.
